% ---
% id: foundations/lem:binary64-lattice
% kind: lem
% title: "Binary64 numbers on dyadic lattices"
% summary: >-
%   Every lattice value at an exponent k ≤ 53 is a binary64 number, and between
%   two consecutive powers of two the binary64 numbers are the multiples of one
%   power of two, in [1/2, 1) those of two to the power −53. Representability of
%   a declared entry is therefore a divisibility condition.
% prominence: standard
% status: accepted
% uses: [def:binary64]
% assumes: []
% refs: []
% provenance: newly-derived
% lean:
% ---

\begin{lemma}\label{lem:binary64-lattice}
\begin{enumerate}[label=(\roman*),leftmargin=2.4em]
\item For every lattice exponent $k$, every real lattice value at $k$ is a
  binary64 number, and every complex lattice value at $k$ lies in
  $\mathbb{B}_{\mathbb{C}}$.
\item For every integer $e$ with $-1022\le e\le1023$,
  \begin{equation}\label{eq:binade}
    \mathbb{B}\cap\big[2^{e},2^{e+1}\big)
    \;=\;2^{e-52}\,\mathbb{Z}\cap\big[2^{e},2^{e+1}\big).
  \end{equation}
  In particular the binary64 numbers in $[\tfrac12,1)$ are the multiples of
  $2^{-53}$ in that interval, and those in $[1,2)$ are the multiples of
  $2^{-52}$ in it.
\end{enumerate}
\end{lemma}

\begin{proof}
(i) Let $x$ be a real lattice value at $k$ (Definition~\ref{def:binary64}\,(iv))
and write $x=n\,2^{-k}$ with $n\in\mathbb{Z}$; then
$|n|=2^{k}|x|\le2^{k}\le2^{53}$. If $|n|<2^{53}$, then $x=n\,2^{q}$ with
$q=-k\in[-53,-1]\subseteq[-1074,971]$, so $x\in\mathbb{B}$ by
Definition~\ref{def:binary64}\,(i). Otherwise $|n|=2^{53}$, which forces
$k=53$ and $|x|=1$, and $x=(\pm1)\,2^{0}\in\mathbb{B}$. A complex lattice value
has its real and imaginary parts on the lattice, hence in $\mathbb{B}$.

(ii) Let $x=n\,2^{e-52}$ with $n\in\mathbb{Z}$ and $2^{e}\le x<2^{e+1}$. Then
$2^{52}\le n<2^{53}$, and $q\coloneqq e-52$ satisfies $-1074\le q\le971$
because $-1022\le e\le1023$; so $x\in\mathbb{B}$. Conversely, let
$x=m\,2^{q}\in\mathbb{B}$, as in Definition~\ref{def:binary64}\,(i), with
$2^{e}\le x<2^{e+1}$. Then $m>0$, and $m<2^{53}$ gives
$2^{q}=x/m>2^{e}\,2^{-53}$, so the integer $q$ satisfies $q>e-53$, that is
$q\ge e-52$. Hence $x=\big(m\,2^{q-e+52}\big)\,2^{e-52}$ with
$m\,2^{q-e+52}\in\mathbb{Z}$, which proves \eqref{eq:binade}. The two
particular cases are $e=-1$ and $e=0$.
\end{proof}
